%% ---------------------------------------------------------------------------
%% DeLP audit paper (v0.2) -- Where the gains come from
%% Target: IJCNN 2027 (IEEEtran conference style).  Compile: pdflatex -> bibtex -> pdflatex x2
%% Figures live in ../figures (300 dpi PNG).
%% ---------------------------------------------------------------------------
\documentclass[conference]{IEEEtran}
\usepackage{iftex}
\ifPDFTeX
  \usepackage[T1]{fontenc}
  \usepackage[utf8]{inputenc}
  \usepackage{newunicodechar}
  \newunicodechar{−}{\ensuremath{-}}
  \newunicodechar{±}{\ensuremath{\pm}}
  \newunicodechar{×}{\ensuremath{\times}}
  \newunicodechar{→}{\ensuremath{\rightarrow}}
  \newunicodechar{≥}{\ensuremath{\geq}}
  \newunicodechar{≤}{\ensuremath{\leq}}
  \newunicodechar{≈}{\ensuremath{\approx}}
  \newunicodechar{—}{---}
  \newunicodechar{–}{--}
  \newunicodechar{’}{'}
  \newunicodechar{“}{``}
  \newunicodechar{”}{''}
  \newunicodechar{ℓ}{\ensuremath{\ell}}
\fi
\usepackage{amsmath,amssymb}
\usepackage{graphicx}
\usepackage{booktabs}
\usepackage{url}
\graphicspath{{./figures/}}

% --- mild spacing tightening so the paper fits the 8-page limit (content unchanged)
\setlength{\textfloatsep}{8pt plus 2pt minus 2pt}
\setlength{\floatsep}{7pt plus 2pt minus 2pt}
\setlength{\intextsep}{7pt plus 2pt minus 2pt}
\setlength{\abovedisplayskip}{4pt plus 1pt minus 1pt}
\setlength{\belowdisplayskip}{4pt plus 1pt minus 1pt}
\setlength{\abovedisplayshortskip}{2pt plus 1pt}
\setlength{\belowdisplayshortskip}{2pt plus 1pt}

\newcommand{\delp}{\textsc{DeLP}}
\newcommand{\bwt}{\mathrm{BWT}}

\begin{document}

\title{Where Do the Gains Come From? An Audit of\\ Dendritic Context-Gated Protection\\ for Continual Learning}

\author{\IEEEauthorblockN{Ranxi Zhou}
\IEEEauthorblockA{School of Information Science and Engineering\\
Hunan University of Chinese Medicine\\
Changsha 410208, China\\
Email: \texttt{zhouranxi@foxmail.com}}}

\maketitle

\begin{abstract}
Cortical pyramidal neurons integrate contextual signals in dendritic compartments, and networks
with active dendrites inherit this context selectivity. In preliminary, unpublished experiments we
proposed \delp{}
(Dendritic Localized Protection), which freezes dendritic segments for non-current contexts and
weights an elastic penalty by each neuron's measured usage under previous tasks; on Permuted-MNIST
it appeared to dominate every baseline in both forgetting and accuracy. This paper is the audit we
should have run before publication. Prompted by an anomaly --- a baseline whose penalty was insensitive to a
100-fold change of its only hyperparameter --- we found two defects in our MAS implementation,
three deviations from the SI reference, a fixed-fraction pruning schedule that exhausted PackNet's
capacity, a missing plain-EWC arm, hyperparameters selected on the test set, and $\lambda$ grids
that left every baseline in its protection-free regime. After correction, \delp{} keeps the best
five-task accuracy (80.37\%, next best 75.27, $p=0.0024$) but no longer leads on forgetting, and at
ten tasks corrected EWC and MAS exceed its accuracy. A seven-rung mechanism ladder localizes what
remains: with baselines given the same hidden activation, the entire dendritic chain adds
$+0.37$ points at five tasks ($p=0.49$); the one mechanism that survives a parameter-budget cap is
usage weighting ($+1.79$, $p=0.005$ at twenty tasks), worth $+4$ to $+6$ points once both sides are
gain-calibrated. A hidden-layer gain scalar alone moves accuracy by 16 points --- more than the
entire method --- so the original comparison measured an uncalibrated gain, not a mechanism.
Finally, preregistered experiments in a rule-switching task matched to the neuroscience show that
compartmentalized routing itself confers no adaptation advantage larger than 30\% --- the
sensitivity limit of 8 preregistered seeds --- across three cue-conflict intensities, up to and
including fully adversarial cues. We report the corrected accounting as the contribution: when
benchmarking protection methods, the gain must be declared, the grid must bracket the optimum, and
the baseline must be implemented as carefully as the method.
\end{abstract}

\begin{IEEEkeywords}
dendritic computation, continual learning, catastrophic forgetting, active dendrites, elastic
weight consolidation, reproducibility
\end{IEEEkeywords}

%% ===========================================================================
\section{Introduction}
%% ===========================================================================
Cortical pyramidal neurons are not point units: their dendritic arbors form semi-independent
compartments that integrate contextual signals with sensory input, and recent in vivo work shows
that apical dendrites encode abstract rules which the soma combines with sensory
evidence~\cite{zhang2026}. Networks with active dendrites~\cite{iyer2022} inherit this
context selectivity, and it is tempting to read the same structure as a recipe for continual
learning: if dendritic gating is the biological solution to interference within a cell, the
structure should indicate \emph{where} an elastic penalty~\cite{kirkpatrick2017,zenke2017} should
protect, and how much~\cite{mallya2018}.

We pursued this idea in preliminary, unpublished experiments (archived; Appendix~A) and obtained
striking support: a usage-localized elastic penalty (\delp{})
appeared to reduce forgetting from $-12.67$ to $-1.94$ BWT and raise accuracy from 74.96\% to
80.37\%, beating Synaptic Intelligence (SI), MAS and PackNet on both axes at five tasks. That
comparison, however, carried defects we had not looked for --- one instance of a pattern repeated
across machine-learning subfields, where re-tuned baselines and honest hyperparameter budgets
erase apparent advances~\cite{musgrave2020,lucic2018,melis2018}. This paper reports the audit, the
corrected results, and two lines of evidence that were only possible after it.

\noindent\textbf{How the audit started.} In our screening data, MAS was insensitive to a
100-fold change of its only hyperparameter ($\lambda\in\{0.1,1,10\}$ moved its BWT by $0.26$
points, versus $8.35$ for SI and $8.50$ for \delp{} on the same protocol). A penalty that does not
respond to its own strength is not being applied. Chasing that one observation exposed two MAS
bugs, three deviations from the SI reference, a capacity-exhausting PackNet schedule, a missing
EWC arm, and two protocol errors (test-set selection and degenerate $\lambda$ grids) that
systematically favored the proposed method.

\noindent\textbf{What survives.} We re-ran the full comparison with every arm implemented and
selected correctly (Section~\ref{sec:audit}); decomposed the method into seven mechanisms on one
pipeline (Section~\ref{sec:ladder}); and ran three preregistered experiments in a rule-switching
task matched to the neuroscience, asking whether compartmentalized routing itself helps adaptation
(Section~\ref{sec:prereg}). Three conclusions survive, all smaller than the original claims:
usage weighting is the one dendritic mechanism with evidence that survives budget capping and gain
calibration; the dominant factor in the original table was an uncalibrated hidden-layer gain; and
compartmental structure, on its own, buys nothing measurable in either paradigm.

\noindent\textbf{Contributions.}
\begin{enumerate}
\item A defect-by-defect audit of our own preliminary comparison, with corrected eight-arm results at
five and ten tasks and paired tests (Section~\ref{sec:audit}).
\item A mechanism ladder isolating seven candidate sources of gain, plus a gain-sweep analysis
showing that a single scalar --- hidden-layer gain --- outweighs the entire method, with a probe
showing the proposed gate never opens (Section~\ref{sec:ladder}).
\item A calibration protocol separating ``default'', ``capped-budget'' and ``gain-calibrated''
comparisons, under which usage weighting is the only mechanism that survives
(Section~\ref{sec:calib}).
\item Preregistered rule-switching experiments (with instrument-incident gates) showing no
adaptation benefit of compartmentalized routing across three cue-conflict intensities, with a post
hoc sensitivity analysis bounding the claim to effects above a 30\% adaptation reduction
(Section~\ref{sec:prereg}).
\end{enumerate}

%% ===========================================================================
\section{Background and Related Work}
%% ===========================================================================
\noindent\textbf{Dendritic computation.} Cable theory and two-layer models made the theoretical
case for dendritic nonlinearity~\cite{poirazi2003}, and its computational vocabulary is mapped in
detail~\cite{london2005,poirazi2020}; one cortical neuron can approximate a deep
network~\cite{beniaguev2021}, and apical dendrites were proposed as excitable integrators long
ago~\cite{chang1951,ujfalussy2018}. Dendritic-plasticity work argues that context-dependent
inhibition gates plasticity at branches, letting overlapping assemblies coexist~\cite{onasch2026}.
On the engineering side, dendritic complexity is argued to be a route to neuromorphic
efficiency~\cite{cardwell2025,boyle2025,baek2024,boahen2022}.

\noindent\textbf{Active dendrites and continual learning.} Neurons with multiple dendritic
segments that receive a context vector and gate the feedforward response support multi-task
learning~\cite{iyer2022}. Regularization methods protect important weights (EWC's
Fisher~\cite{kirkpatrick2017}, SI's path integral~\cite{zenke2017}, MAS's output-norm
sensitivity~\cite{aljundi2018}); isolation methods allocate capacity (PackNet~\cite{mallya2018},
hard attention~\cite{serra2018}); context-dependent gating modulates units by
context~\cite{masse2018}. \delp{} sits between the families: a Fisher-style penalty whose
per-parameter weight comes from the gating structure itself~\cite{parisi2019}. The audit protocol
used here --- validation-selected grids, parameter-budget caps, gain calibration, preregistered
gates --- follows the reporting standards advocated for empirical
ML~\cite{lipton2019,agarwal2021}.

%% ===========================================================================
\section{Method}\label{sec:method}
%% ===========================================================================
\begin{figure*}[t]
\centering
\includegraphics[width=\textwidth]{fig_concept.png}
\caption{(a) The active-dendrite neuron: a context vector is projected onto $K$ dendritic
segments, the strongest segment gates the shared feedforward drive, and a kWTA nonlinearity
sparsifies the layer; only the winning segment is updated during a task. (b) The three candidate
interference channels and their measured consequences at five tasks: freezing the dendritic
segments changes little; per-task readouts relocate the interference; hard per-neuron freezing
collapses capacity.}
\label{fig:concept}
\end{figure*}

\noindent\textbf{Active-dendrite neuron.} Each hidden neuron $j$ ($j=1..N$) carries $K$ dendritic
segments. A context vector $\mathbf{c}\in\mathbb{R}^{T}$ (one-hot task identity) is projected onto
segments through $U\in\mathbb{R}^{K\times N\times T}$, the dendritic response is the maximum over
segments, and
\begin{equation}
h_j = \mathrm{kWTA}\big( (\mathbf{w}_j^{\top}\mathbf{x} + b_j)\cdot \sigma(d_j) \big),
\end{equation}
with $\sigma$ the logistic function and kWTA keeping the top-$k$ activations per layer
($k=N/4$). During training of task $t$, gradients to $U$ are masked so only the rows selected by
context $t$ are updated (segment freeze).

\noindent\textbf{Why the penalty looks at the shared projection.} An interference-channel
ablation established where forgetting lives: freezing dendritic segments changes little
($-14.79$ vs.\ $-15.16$ BWT for unfrozen); extending the freeze to the feedforward rows of
task-owned neurons collapses accuracy to 67.35\% (capacity starved); per-task readouts relocate the
interference into the shared feedforward projection ($-19.78$ BWT). The shared feedforward
projection and readout are therefore the channels a protection method must address, and the
dendritic gate supplies a per-neuron \emph{usage} statistic for grading the protection:
$o_{t',j}$, the fraction of task-$t'$ samples on which neuron $j$ is inside the kWTA winner set.
The penalty is
\begin{equation}
\mathcal{L} = \mathcal{L}_{\mathrm{CE}} + \lambda \sum_{t' < t} \sum_{i} F^{\mathrm{loc}}_{t',i}\,
(\theta_i - \theta^{\ast}_{t',i})^2, \quad
F^{\mathrm{loc}}_{t',i} = o_{t',j(i)} \cdot F_{t',i},
\end{equation}
with $F_{t'}$ the diagonal Fisher estimated on 300 samples per task and $o$ on 500; both reuse the
task's training data. Unlike hard gating, protection is soft: a penalized neuron can still move at
a quadratic price. Fisher uses the max-log-probability variant throughout (Section~\ref{sec:lam}).
Test-time context inference from a 500-sample buffer (linear classifier on feedforward features)
matched oracle context exactly.

%% ===========================================================================
\section{Auditing the original comparison}\label{sec:audit}
%% ===========================================================================
The original comparison (five seeds, paired $t$-tests, one pipeline) reported \delp{} at
$-1.94\pm0.42$ BWT / $80.37\pm0.67$\% accuracy, ahead of SI ($-3.39$ / 77.56), MAS ($-12.32$ /
75.27) and PackNet (0.00 / 43.95) --- these are the archived original numbers
(\texttt{exp12\_result.json}); the SI rows of Table~\ref{tab:main} report the corrected re-run
under the validation-selected protocol and therefore differ. The audit found seven defects; we
state each with its location and consequence.

\noindent\textbf{D1: MAS importance estimated from batch-averaged gradients, twice divided.} Our
MAS backward-passed an already-averaged objective, multiplied squared \emph{mean} gradients by
$1/B$ again (losing sample variance and shrinking $\Omega$ by the batch size), and then divided the
accumulated $\Omega$ by the usage count at the end of \emph{every} task without resetting, so task
0's contribution decayed as $(1/8)^{t}$ ($7.5\times10^{-9}$ by task 9). The penalty protected only
the most recent task. Consequence: MAS's $\lambda$ elasticity was $+0.08$ over two orders of
magnitude; after the fix it is $+1.40$, and its best five-task BWT improves from $-12.07$ to
$-3.60$. The original ``MAS trails by 10.4 points'' was mostly this bug.

\noindent\textbf{D2: SI deviated from its reference in three places.} The denominator used
$|\Delta\theta|$ instead of $(\Delta\theta)^2+\xi^2$ (with $\xi=10^{-3}$ instead of $0.1$), and an
\texttt{abs()} destroyed the sign of the path integral. The effect on BWT at matched optimum is
small ($-0.15$ vs.\ $-0.36$), but the implementation was not SI.

\noindent\textbf{D3: PackNet pruned 20\% of free capacity per task.} After five tasks 32.8\% of
parameters remain, after ten 10.7\%: the reported collapses to 43.95\% / 27.38\% are budget
exhaustion, not method failure. The corrected variant plans per-task capacity proportionally and
prunes per layer.

\noindent\textbf{D4: no plain-EWC arm.} \delp{} \emph{is} an EWC penalty plus branch protection, so
without a plain-EWC arm the decomposition ``how much is EWC, how much is the dendritic part'' was
unanswerable. Added.

\noindent\textbf{D5: hyperparameters selected on the test set.} Screening evaluated $\lambda$
candidates on the test set. Corrected protocol: held-out validation split (train[10000:12000]).

\noindent\textbf{D6: degenerate $\lambda$ grids.} The best $\lambda$ of \delp{}, SI and MAS all sat
on the upper bound of their grids: the baselines were left in their protection-free regime while
the proposed method's grid covered its working range. Extending SI's grid to $\lambda=100$ moves
its BWT from $-4.04$ to $-0.36$. A related selection rule --- choosing $\lambda$ by best BWT ---
is structurally degenerate here, since BWT improves monotonically with $\lambda$ while accuracy is
unimodal: across 21 (mechanism $\times$ task-count) cells, BWT-based selection hit the grid upper
bound 21 times, each 8--31 accuracy points below the cell's peak. All corrected results select
$\lambda$ by validation accuracy.

\noindent\textbf{D7: initialization not seeded (protocol).} An earlier unseeded pipeline drifted in
the third digit and once flipped a verdict (SI advantage $p=0.058$ unseeded vs.\ $p=0.039$
seeded). All corrected runs seed every arm.

\noindent\textbf{Withdrawn suspicion.} Architectural parameter inequality between the proposed
arm and the baselines ($+1.25$\%) was re-measured and is negligible; we withdraw it as a concern.

\begin{table*}[t]
\centering
\caption{Corrected comparison, five seeds, paired $t$-tests against \delp{}. Every arm uses its
own best $\lambda$ selected on held-out validation (grids per family are released with the code
and verified to bracket each optimum, \texttt{grid\_check.md}); \emph{orig.} rows reproduce the
original implementations (verified bit-identical on shared seeds). BWT: higher is better; acc:
higher is better.}
\label{tab:main}
\small
\begin{tabular}{lrrrr}
\toprule
 & \multicolumn{2}{c}{Five tasks} & \multicolumn{2}{c}{Ten tasks} \\
\cmidrule(lr){2-3}\cmidrule(lr){4-5}
Method & BWT & Final avg acc & BWT & Final avg acc \\
\midrule
\delp{} (ours)                     & $\mathbf{-1.94 \pm 0.42}$ & $\mathbf{80.37 \pm 0.67}$ & $-5.38 \pm 0.44$ & $55.23 \pm 1.41$ \\
EWC (online, multi-anchor)         & $-7.08 \pm 1.19$ & $74.28 \pm 0.86$ & $-7.81 \pm 1.28$ & $\mathbf{68.35 \pm 0.95}$ \\
SI (reference impl.)               & $\mathbf{-0.03 \pm 0.35}$ & $62.33 \pm 1.41$ & $\mathbf{-1.31 \pm 0.29}$ & $40.57 \pm 1.51$ \\
SI (orig.\ impl.)                  & $-0.15 \pm 0.38$ & $63.34 \pm 1.63$ & $-0.86 \pm 0.13$ & $39.46 \pm 1.46$ \\
MAS (corrected)                    & $-3.72 \pm 0.82$ & $75.22 \pm 0.55$ & $-5.22 \pm 0.66$ & $62.91 \pm 1.02$ \\
MAS (orig.\ impl.)                 & $-12.32 \pm 2.53$ & $75.27 \pm 1.99$ & $-19.88 \pm 1.27$ & $68.00 \pm 1.23$ \\
PackNet (corrected)                & $-12.67 \pm 2.60$ & $74.96 \pm 2.01$ & $-19.92 \pm 0.68$ & $67.93 \pm 0.83$ \\
PackNet (orig.\ impl.)             & $0.00$ (mask) & $43.95 \pm 0.82$ & $0.00$ (mask) & $27.38 \pm 0.63$ \\
\bottomrule
\end{tabular}
\end{table*}

\noindent\textbf{Corrected comparison.} Table~\ref{tab:main} replaces the original main table.
Three original claims do not survive. (i) ``MAS trails by 10 points'': the corrected MAS trails by
1.8 accuracy-\emph{significant} points at five tasks ($p=0.015$ on BWT, $p=0.0001$ on accuracy) and
is statistically indistinguishable on BWT at ten tasks ($-5.22$ vs.\ $-5.38$, $p=0.63$) with
7.7 points \emph{more} accuracy ($p=0.0001$). (ii) ``\delp{} forgets less than SI'': at ten tasks
SI's BWT is significantly better ($-1.31$ vs.\ $-5.38$, $p<0.0001$). (iii) ``best accuracy at ten
tasks'': corrected EWC (68.35\%, $p=0.0003$) and MAS (62.91\%, $p=0.0001$) both exceed \delp{}'s
55.23\%. What survives: the best five-task accuracy, 5.1 points above the next best arm
($p=0.0024$); the orig.\ rows show the original table compared against broken baselines. The table
is best read as a frontier, not a ranking: SI occupies the low-forgetting end, corrected EWC/MAS
the high-accuracy end at ten tasks, and \delp{} the best measured five-task operating point.

\begin{table}[t]
\centering
\caption{Corrected comparison on Permuted-Fashion-MNIST (second benchmark; 5 seeds, val-selected
$\lambda$, same protocol as Table~\ref{tab:main}). Original-implementation rows, omitted here,
again trail their corrected counterparts by 5--11 accuracy points.}
\label{tab:fashion}
\scriptsize
\setlength{\tabcolsep}{3.5pt}
\begin{tabular}{@{}lrrrr@{}}
\toprule
 & \multicolumn{2}{c}{Five tasks} & \multicolumn{2}{c}{Ten tasks} \\
\cmidrule(lr){2-3}\cmidrule(l){4-5}
Arm & acc. & BWT & acc. & BWT \\
\midrule
\delp{}           & \textbf{66.9}$\pm$0.9 & $-8.5$$\pm$1.4 & 55.5$\pm$1.6 & $-18.0$$\pm$1.6 \\
MAS (corrected)   & 60.8$\pm$0.7 & $-4.1$$\pm$1.5 & \textbf{56.3}$\pm$1.1 & $-8.5$$\pm$1.3 \\
EWC               & 57.1$\pm$1.8 & $-11.3$$\pm$1.2 & 53.5$\pm$3.4 & $-12.7$$\pm$3.0 \\
PackNet (corr.)   & 55.3$\pm$3.1 & $-15.2$$\pm$3.2 & 48.4$\pm$3.9 & $-22.2$$\pm$4.3 \\
SI (corrected)    & 34.6$\pm$2.7 & $\mathbf{-2.8}$$\pm$0.9 & 26.6$\pm$1.2 & $\mathbf{-1.9}$$\pm$0.6 \\
\bottomrule
\end{tabular}
\end{table}

\noindent\textbf{Second benchmark.} Does the corrected picture transfer? Re-running the full
corrected protocol (identical arms, grids, val $\lambda$ rule and seeds) on Permuted-Fashion-MNIST
gives the same shape (Table~\ref{tab:fashion}): \delp{} holds the best five-task accuracy
($+6.2$ over corrected MAS, $p=0.0003$), the ten-task accuracy lead disappears (MAS $+0.8$,
$p=0.54$; EWC $p=0.32$), and corrected SI/MAS forget significantly less at ten tasks
($p\leq0.003$). The conclusion of the audit is not an artifact of one dataset.

%% ===========================================================================
\section{Mechanism ladder and the gain confound}\label{sec:ladder}
%% ===========================================================================
The corrected table says the \emph{comparison} was flawed, not which part of the method works. To
attribute gains to mechanisms we trained seven rungs on one pipeline, each rung adding exactly one
mechanism, all else identical (R0: point MLP + EWC with squashed hidden activation;
R0b: value-range compression removed; R1: $+$ dendritic parameters, context pinned to zero;
R2: $+$ context gating; R3: $+$ segment freeze; R4: $+$ usage weighting ($=$ \delp{}); R5: $+$ CE
Fisher). Two self-checks gate the analysis: R1 must be bit-identical to R0b (it is: $\Delta=0.00$,
$p=1.000$, confirming the padded dendritic layer is dead weight), and every rung's $\lambda$
elasticity must be positive (all are, $+2.55$ to $+7.94$ in the released ladder batch).

\begin{table}[t]
\centering
\caption{Rung-by-rung accuracy gain (each rung at its own validation-selected $\lambda$; $p$ in
parentheses; five seeds). Differences are computed from unrounded per-seed means, so
subtracting the rounded rung accuracies in the released table can differ by $0.01$.}
\label{tab:ladder}
\footnotesize
\setlength{\tabcolsep}{3.5pt}
\begin{tabular}{@{}lrrr@{}}
\toprule
Step & $T{=}5$ & $T{=}10$ & $T{=}20$ \\
\midrule
R0$\to$R0b: remove squashing   & $+5.37$ (.000) & $+4.49$ (.000) & $+2.13$ (.012) \\
R0b$\to$R1: dendritic params   & $+0.00$ (1.00) & $+0.00$ (1.00) & $+0.00$ (1.00) \\
R1$\to$R2: context gating      & $+0.61$ (.172) & $+0.37$ (.019) & $+1.35$ (.260) \\
R2$\to$R3: segment freeze      & $-0.11$ (.616) & $+0.05$ (.873) & $-0.37$ (.602) \\
R3$\to$R4: usage weighting     & $-0.13$ (.563) & $+0.07$ (.820) & $+1.02$ (.008) \\
R4$\to$R5: CE Fisher           & $+0.04$ (.842) & $+0.04$ (.863) & $-0.76$ (.404) \\
\midrule
R4 (full method) vs.\ R0b      & $+0.37$ (.494) & $+0.49$ (.083) & $+2.00$ (.025) \\
\bottomrule
\end{tabular}
\end{table}

Table~\ref{tab:ladder} is the paper's central accounting. The single largest step is the removal
of the hidden squashing: sigmoid compresses the pre-activation into $(0,1)$, collapsing amplitude
information (a gain sweep shows accuracy rises monotonically with the post-sigmoid gain,
$r=+0.94$ --- amplitude compression, not ``saturation'', is the mechanism). The dendritic
parameters themselves contribute exactly zero (bitwise identity). Context gating is small, positive
and direction-stable (BWT $+3.10$, $p=0.001$ at five tasks). Segment freeze is zero. Usage
weighting is the only rung whose value grows with sequence length ($+1.02$, $p=0.008$ at twenty
tasks), as expected if it offsets penalty amortization across more anchors. Summed against the
same-activation baseline R0b, the full method is $+0.37$ ($p=0.49$) at five tasks, $+0.49$
($p=0.083$) at ten, and $+2.00$ ($p=0.025$) at twenty.

\noindent\textbf{The gain confound, quantified.} kWTA is scale-equivariant
($\mathrm{kWTA}(\alpha t)=\alpha\,\mathrm{kWTA}(t)$, verified exactly), so a hidden-layer gain
$\alpha$ does not change the function class --- but it changes how far weights must travel, and
EWC-style penalties are quadratic in displacement. Sweeping $\alpha$ on plain ``MLP+EWC'' moves
five-task accuracy from 76.24 ($\alpha_{\mathrm{eff}}\approx0.5$) to 92.40 ($\alpha{=}16$) ---
$+16.2$ points, versus $+0.37$ for the entire method chain; at $\alpha{=}2$ already, plain
MLP+EWC beats the full \delp{} on accuracy \emph{and} forgetting at all three task counts. The
\delp{} gate is $t\cdot\sigma(d)$ with $d$ learnable, so it could in principle open; a probe that
logs the gate during full training shows it never does: $\overline{\sigma(d)}=0.575$--$0.577$
across all five tasks (0.55 at twenty), i.e., an uncalibrated $\approx0.58$ gain sitting near the
worst end of the sweep. The original comparison therefore measured a gain miscalibration, not a
mechanism.

\subsection{What survives: calibration protocols}\label{sec:calib}
Because gain interacts with every $\lambda$-style penalty, ``how much better is the method''
depends on the calibration protocol, and we now report all three (Table~\ref{tab:proto}).
Under \emph{default} settings (neither side gain-tuned; dendritic parameter budget capped at
$nseg=\min(8,T)$, which cuts the twenty-task dendritic block from 102{,}400 to 40{,}960 parameters
with all five rungs unchanged within noise), usage weighting beats the same-activation baseline by
$+1.79$ ($p=0.005$) at twenty tasks. Under \emph{gain calibration} (two-point gain probe per
family at each family's own optimal $\lambda$), \delp{} gains $+6.3$/$+8.6$/$+9.7$ points
(five/ten/twenty tasks) --- but so do the baselines (MAS $+4.1$/$+3.4$/$+4.6$, SI
$+3.3$/$+6.0$/$+3.1$, EWC $+3.5$/$+0.7$/$+3.0$); the calibrated comparison leaves \delp{} ahead by
$+4$ to $+6$ points (87.81/83.00/75.18 vs.\ best baseline 81.80/78.99/70.92). Under the
\emph{original} protocol (neither side tuned, published configurations) the difference
\emph{changes sign} with $T$ ($-0.85$ to $+3.7$) and is unusable. The usage-weighting result is a
per-family two-point probe at fixed $\lambda^{\ast}$, not a joint (gain $\times$ $\lambda$) grid,
so $+4$--$6$ is an approximate lower bound.

\begin{table}[t]
\centering
\caption{The same question under three calibration protocols: \delp{} accuracy advantage over the
best same-activation baseline.}
\label{tab:proto}
\footnotesize
\setlength{\tabcolsep}{3pt}
\begin{tabular}{@{}p{0.36\columnwidth}rp{0.33\columnwidth}@{}}
\toprule
Protocol & Adv. & Status \\
\midrule
Both default (published cfg.) & $-0.85$ to $+3.7$ & changes sign with $T$; unusable \\
Usage weighting, capped budget & $+1.79$ ($p{=}.005$) & $T{=}20$ only; budget-controlled \\
Both gain-calibrated & $+4$ to $+6$ & two-point probe; lower bound \\
\bottomrule
\end{tabular}
\end{table}

%% ===========================================================================
\section{Preregistered rule-switching tests}\label{sec:prereg}
%% ===========================================================================
The biological motivation predicts more than a usage statistic: if apical compartments are where
top-down rule signals meet sensory evidence~\cite{zhang2026}, a network that routes the two into
separate compartments should \emph{adapt faster} when rules change than a network that fuses them.
This is a structural claim, distinct from protection, and we tested it in a preregistered
companion study (thresholds, gates and sample sizes fixed before running; deviations require a
numbered amendment).

\noindent\textbf{Task and models.} Rule-switching categorization: $d=16$ stimuli, 4 random linear
rules (resampled per seed), unsignaled switches every 80 trials, 480-trial episodes, reward
correctness $p=0.85$. GRU-backed policy networks trained with REINFORCE (1500 episodes, 8 seeds,
lr $10^{-3}$, belief-auxiliary loss). DC-Net routes a rule-cue plus history summary and reward
prediction error to a top (apical) compartment feeding a gating signal, while sensory input
reaches the basal dendrites; Point-Net receives the identical information fused into one stream;
Random-MLP has the same parameters with scrambled structure; a trivial-learner gates task
difficulty. In PR-002 each block carries a one-hot rule cue, valid with $p=0.9$ and actively
misleading otherwise, giving blocks where the top-down signal conflicts with sensory evidence ---
the regime where compartmental isolation should help most. PR-003 removes the contrast entirely:
the cue is misleading in \emph{every} block ($p_{\text{valid}}=0$), so the top-down stream carries
pure inverted information and must be suppressed rather than occasionally corrected.

\begin{table}[t]
\centering
\caption{Preregistered rule-switching results across three cue-conflict intensities (cue
misleading in 0\%, 10\%, or 100\% of blocks; per-seed medians of post-switch trials to 80\%
correct, 8 seeds; Mann-Whitney $U$, one-sided; reduction threshold $\geq$25\%).}
\label{tab:prereg}
\scriptsize
\setlength{\tabcolsep}{2.5pt}
\begin{tabular}{@{}lrrr@{}}
\toprule
 & PR-001 (none) & PR-002 (10\%) & PR-003 (100\%) \\
\midrule
DC-Net final acc          & 0.648 & 0.649 & 0.647 \\
Point-Net final acc       & 0.643 & 0.654 & 0.632 \\
Trivial learner acc       & 0.500 & 0.500 & 0.500 \\
DC-Net adapt med.\ [CI]   & 6.5 [4.5, 7.8] & 6.5 [3.9, 7.6] & 6.5 [4.5, 7.8] \\
Point-Net adapt med.      & 7.0 & 7.0 & 6.5 \\
Median reduction          & 7.1\% & 7.1\% & 0.0\% \\
$p$                       & 0.335 & 0.296 & 0.458 \\
Cohen's $d$               & $-0.27$ & $-0.18$ & $-0.37$ \\
\bottomrule
\end{tabular}
\end{table}

\noindent\textbf{Result.} All three hypotheses fail to meet the preregistered criterion
(Table~\ref{tab:prereg}). In PR-001, with validity gates passed (trivial learner 14.8 points below;
scrambled-structure control indistinguishable from DC-Net, so the null cannot be attributed to
parameter count), DC-Net's post-switch adaptation is 7.1\% faster at $p=0.335$. In PR-002, with the
misalignment rate verified at 9.83\% ($n=240{,}000$ switches) against the preregistered 10\%, the
numbers are essentially identical (7.1\%, $p=0.296$); in cue-valid blocks the two architectures are
exactly tied (6.0 vs.\ 6.0, $p=0.69$). In PR-003, with misalignment verified at 100\%, the two are
again exactly tied (6.5 vs.\ 6.5, $p=0.458$) --- and both medians sit at this task's measurement
floor, where detection is hardest. Validity gates passed in all three runs; a first instrumented
run of PR-002 was automatically voided by its own preregistered gate when a logging defect left
condition labels unrecorded, and re-run --- we report this because the gate doing its job is part
of the result.

\noindent\textbf{Sensitivity.} The claims above are bounded, and we bound them ourselves. A post
hoc simulation on the Point-Net reference distribution (20{,}000 resamples per point; DC-Net
samples drawn as a fixed scaling of bootstrap draws) shows that with 8 seeds and a one-sided
Mann-Whitney test at $\alpha=0.05$, the design had 80\% power only against median reductions of
$\geq$30\%; at the preregistered 25\% threshold power was $\approx$50\%, and at the observed 7.1\%
it was $<$5\%. The defensible conclusion is therefore: \emph{no compartmentalization advantage of
30\% or more exists in this task family}, while smaller effects --- including the observed 7.1\% ---
cannot be distinguished from zero with this sample size. In PR-003 both medians sit at the
measurement floor, degrading the detectable effect to 50\%; an exact tie at the floor nonetheless
weakly favors the null.

\noindent\textbf{Reading.} In this task family a GRU history summary already carries the rule
inference; routing top-down information through a separate compartment adds nothing measurable,
at any of three conflict intensities --- none, occasional, or fully adversarial --- and the
conflict gradient itself shows no dose--response trend (adaptation may sit near the
information-theoretic floor for 4-way
rules under 15\% feedback noise). Combined with the ladder: the measurable
value of dendritic structure in our systems is not its routing or its isolation but the
\emph{measurement} it computes --- per-neuron, per-context usage --- and even that value appears
only as a weighting inside an elastic penalty, only under a capped budget, and only at long
horizons.

%% ===========================================================================
\section{Learning versus retention at ten tasks}\label{sec:lam}
%% ===========================================================================
For completeness we retain the decomposition $\bar{a}_T = \ell + \frac{T-1}{T}\bwt$ ($\ell$: mean
just-learned accuracy). At ten tasks the point MLP learns each task to $\approx$86\% and forgets
heavily, while \delp{} learns less and forgets little; the deficit is not a width limit
(256$\to$1024 recovers 3.4 points of $\ell$, costs 2.5 of BWT) nor a representational wall
(single-task ceiling 89.4\% at $N{=}256$), but an accumulated elastic anchor cost: sweeping
$\lambda$ trades the two terms monotonically (BWT improves 16.3 points while $\ell$ falls 33.2
from $\lambda{=}100$ to $10^4$), and truncating the anchor set to the five most recent tasks
recovers part of $\ell$ at every $\lambda$. Per-task kWTA winner sets overlap substantially (mean
pairwise Jaccard 0.51): gating does not partition neurons --- ``ownership'' of parameters is
imposed, not discovered. No single $\lambda$ dominates SI on both axes at ten tasks
(Fig.~\ref{fig:cap}).

\begin{figure}[t]
\centering
\includegraphics[width=0.88\columnwidth]{fig_capacity.png}
\caption{Ten-task learning/retention decomposition: the loss lies in the learning term.}
\label{fig:cap}
\end{figure}

%% ===========================================================================
\section{Discussion}
%% ===========================================================================
\noindent\textbf{What the dendritic story buys --- revised.} The audited evidence supports a
claim with three qualifiers. Dendritic compartmentalization provides no capacity (bitwise-null
when context is pinned), no protection by itself (segment freeze $\approx 0$), and no adaptation
advantage above the 30\% sensitivity limit in preregistered rule-switching tasks across three
conflict intensities, including fully adversarial top-down signals. What it provides, once, and
only as a component: context-conditioned gating computes a per-neuron usage statistic, and
weighting an elastic penalty by that statistic is the one mechanism whose gain survives budget
capping ($+1.79$, $p=0.005$ at twenty tasks) and grows under calibration. This is closer to an
accounting principle than to an architecture: the biology says where to look, standard elastic
machinery does the work, and every other presumed benefit dissolved under controlled comparison.

\noindent\textbf{Lessons for benchmarking protection methods.} Five practices would have caught
most of Section~\ref{sec:audit} before publication: (i) a hyperparameter-elasticity check per arm
(a penalty that ignores its own $\lambda$ is not applied); (ii) validation-based model selection
with grids verified to bracket each arm's optimum on \emph{both} ends; (iii) a plain-EWC arm
whenever the method contains EWC as a component; (iv) an activation- and gain-matched baseline,
since quadratic penalties interact with parameter scale and a single free scalar outweighed the
entire method here; (v) preregistered gates for confirmatory claims, including gates that void the
run when instrumentation fails.

\noindent\textbf{Limitations.} (i) Both benchmarks are MNIST-derived (Permuted-MNIST,
Permuted-Fashion-MNIST); a non-MNIST family such as Split-CIFAR10 remains future
work. (ii) Task identity is available at training time (test-time inference shown lossless here).
(iii) Gain calibration uses a two-point probe per family at fixed $\lambda^{\ast}$, not a joint
(gain $\times$ $\lambda$) grid: the $+4$--$6$ calibrated advantage is an approximate lower bound.
(iv) The one-hot context leaves the capped dendritic block $+20$\% parameters at twenty tasks.
(v) The rule-switching null covers one task family (4-way rules, 85\% feedback noise, GRU
summary); gates cannot rule out effects below $\sim$7\% at this sample size. (vi) All runs use 2
epochs and 10k samples per task.

%% ===========================================================================
\section{Conclusion}
%% ===========================================================================
We asked whether dendritic context gating tells a continual-learning system where to protect its
knowledge, and the audited answer is more modest and more useful than our first report. The
original comparison was inflated by baseline defects, test-set selection and an uncalibrated gain;
after correction, the dendritic structure contributes exactly one mechanism with standing evidence
--- a per-neuron usage statistic that localizes an elastic penalty ($+1.79$ points at twenty tasks
under a capped budget, $+4$--$6$ gain-calibrated) --- and nothing else:
no capacity, no protection from freezing, no adaptation advantage from compartmental routing
across three preregistered conflict levels. A hidden-layer gain scalar moved accuracy more than
the entire method --- both the reason the original table looked as it did and the paper's most
transferable finding. Dendritic computation, in this benchmark family, is most useful not as extra
capacity or isolation but as \emph{measurement}: a signal about which parts of a network mean
what --- provided the comparison is calibrated well enough to see it.

\appendices
\section{Reproducibility and hyperparameters}
Code and raw result files (preregistrations, per-seed JSONs, analysis scripts) are archived at
\url{https://github.com/zhouranxi/delp-audit}. Original experiments:
\texttt{dendritic-lab/} --- \texttt{exp4*}/\texttt{exp5} (ablation chain, \delp{}),
\texttt{exp6} (original baselines), \texttt{exp8} (shared pool, context inference),
\texttt{exp9} (single-pipeline table), \texttt{exp10} (decomposition), \texttt{exp11} (frontier),
\texttt{exp12} (seeded re-run). Audit and re-analysis: \texttt{bench/} ---
\texttt{exp6\_fixed.py} (corrected baselines, validation $\lambda$, EWC arm),
\texttt{ladder.py} (945 cells, zero failures), \texttt{kwta\_scale\_check.py},
\texttt{gate\_probe.py}, \texttt{sig\_probe.py}, \texttt{analyze\_ei.py}. Rule-switching study:
\texttt{dendrite-mve/} with preregistrations \texttt{prereg\_pr00\{1,2,3\}\_*.md} (amendments,
corrections, result records) and raw per-seed JSONs (\texttt{results/}, \texttt{results\_cue/},
\texttt{results\_cue\_adv/}). Every run \texttt{torch.manual\_seed}s before model construction;
settings not stated in the text: $N=256$, $K=T$ segments ($8$, top-3, in the capped variant),
kWTA $k=N/4$, Fisher/usage from 300/500 samples, five seeds (eight for the rule-switching study),
one RTX 3090. Corrected-baseline runs total 80 cells; ladder and calibration runs 665 cells.

%% ===========================================================================
\section*{Acknowledgment}
%% ===========================================================================
The author used large language model tools (DeepSeek V4.1 and GLM) to draft and edit the English
prose and to implement and verify the audit scripts under the author's direction. The experimental
design, the defect diagnosis, and all scientific claims were made by the author, who takes full
responsibility for the content of this article.

\bibliographystyle{IEEEtran}
\bibliography{refs}

\end{document}
